\documentclass[11pt,a4paper]{article}
\usepackage[utf8]{inputenc}
\usepackage[T1]{fontenc}
\usepackage{amsfonts}
\usepackage[french]{babel}
\frenchbsetup{StandardLists=true}
\usepackage{enumitem}
\usepackage{amssymb}
\def\nbR{\ensuremath{\mathrm{I\! R}}}
\usepackage{amsmath}
\usepackage{awesomebox}
\DeclareMathOperator{\e}{e}

\usepackage[left=2cm,right=2cm,top=2cm,bottom=2cm]{geometry}
\usepackage{multirow}
\usepackage[dvipsnames]{xcolor}
\usepackage{parskip}
\usepackage{caption}

\usepackage[T1]{fontenc}
\usepackage{fouriernc}
\usepackage{graphicx}
\usepackage{setspace}
\singlespacing
\usepackage{tcolorbox}
\usepackage{framed}
\usepackage{multicol}

\usepackage{fancybox}
\usepackage{fancyhdr}
\setlength{\headheight}{14pt}
\pagestyle{fancy}
\renewcommand\headrulewidth{1pt}
\fancyhead[L]{Lycée Denis Diderot}
\fancyhead[R]{\textcolor{red}{\textbf{VERSION PROFESSEUR}} -- 1BTS}
\setcounter{page}{1}\fancyfoot[C]{page \thepage}
\fancyfoot[R]{\today}
\fancyfoot[L]{Alexis RUIZ}

\usepackage{tikz}
\usepackage{pgfplots}
\pgfplotsset{compat=1.18}
\usepackage{tkz-tab}
\usetikzlibrary{arrows}
\usepackage{pifont}

\usepackage{multirow,tabularx,booktabs}
\usepackage{colortbl}
\usepackage[np]{numprint}
\usepackage[tikz]{bclogo}

\def\N{{\mathbb N}}
\def\Z{{\mathbb Z}}
\def\R{{\mathbb R}}
\def\C{{\mathbb C}}
\def\e{{\text{e}}}
\def\dt{{\text{d}t}}

\newcommand{\ud}{\mathrm{d}}

% Boîte de correction bleue
\newcommand{\cor}[1]{%
  \begin{tcolorbox}[
    colback=blue!5!white,
    colframe=blue!70!black,
    boxrule=1pt,
    arc=3pt,
    left=4pt, right=4pt, top=3pt, bottom=3pt,
    before skip=2mm, after skip=3mm
  ]
  \color{blue!80!black}#1
  \end{tcolorbox}%
}

% Boîte de résultat (vert)
\newcommand{\res}[1]{%
  \begin{tcolorbox}[
    colback=green!5!white,
    colframe=green!50!black,
    boxrule=0.8pt,
    arc=3pt,
    left=4pt, right=4pt, top=2pt, bottom=2pt,
    before skip=1mm, after skip=2mm
  ]
  \color{green!40!black}#1
  \end{tcolorbox}%
}

% Compteur et boîte pour les exercices
\newcounter{numexos}
\setcounter{numexos}{0}
\newcommand{\bex}[2]{%
\def\bextitle{#1}%
\addtocounter{numexos}{1}%
\begin{tcolorbox}[colback=white,colframe=black!30,arc=3pt,boxrule=0.6pt,
  left=4mm,right=4mm,top=1mm,bottom=2mm]%
  \textbf{Exercice~\thenumexos\ifx\bextitle\empty\else~:~~#1\fi}\par
  #2%
\end{tcolorbox}}

% Compteur et boîte pour les propriétés
\newcounter{nump}
\setcounter{nump}{0}
\newcommand{\bprop}[2]{%
\addtocounter{nump}{1}%
\begin{tcolorbox}[colback=white,colframe=teal!60,arc=3pt,boxrule=0.8pt,
  left=4mm,right=4mm,top=1mm,bottom=2mm]%
  \textbf{Propriété~\thenump\ifx\relax#1\relax\else~:~~#1\fi}\par
  \smallskip #2%
\end{tcolorbox}}

% Réponse inline en bleu
\newcommand{\ans}[1]{\textcolor{blue!80!black}{\textbf{#1}}}

\newcommand{\cadreblanc}[1]{%
\medskip\framebox[\linewidth][c]{%
\rule{0mm}{#1mm}}\par
\medskip}

\setlength{\columnseprule}{.25mm}
\newcolumntype{C}[1]{>{\centering\arraybackslash}p{#1cm}}
\renewcommand{\arraystretch}{1.4}

\newtcolorbox{mybox}{colback=white!10!white,colframe=black!5!black}

\newcommand{\bdef}[1]
{\begin{bclogo}[couleur = white,logo=\bcinfo , arrondi = 0.1]
{~Définition~:}
~~
{#1}
~~
\end{bclogo}}

\newcommand{\rem}[1]
{\begin{bclogo}[couleur = white,logo=\bcinfo , arrondi = 0.1]
{~Remarque~:}
~~
{#1}
~~
\end{bclogo}}

\newcommand{\bpre}[1]
{\begin{bclogo}[couleur = white,logo=, barre=none, arrondi = 0.1]
{~Prérequis~:}
~~
{#1}
~~
\end{bclogo}}

%%%%%%%%%%%%%%%%%%%%%%%%%%%%%%%%%%%%%%%%%%%%%%%%%%%%%%%%%%%%%
\begin{document}

\begin{mybox}
\begin{center}
\textbf{\LARGE CHAPITRE 6 -- Probabilités conditionnelles}
\end{center}
\end{mybox}

\vfill

\begin{center}
\begin{minipage}{14cm}
\begin{bclogo}[couleur = white, arrondi = 0.1]
{~Objectifs~:}
~~
\begin{itemize}
  \item Construire un arbre ou un tableau de probabilités en lien avec une situation donnée.
  \item Explorer l'arbre ou le tableau de probabilités pour déterminer des probabilités.
  \item Calculer la probabilité d'un événement connaissant ses probabilités conditionnelles relatives à une partition de l'univers.
  \item Utiliser ou justifier l'indépendance de deux événements.
\end{itemize}
\end{bclogo}
\end{minipage}
\end{center}

\vfill

%------------------------------------------------------------------
\bex{Rappels}{}

\begin{enumerate}
  \item $P(A)=0{,}7$, $P(B)=0{,}5$, $P(A\cap B)=0{,}3$. Calculer $P(A\cup B)$.

  \cor{$P(A\cup B)=P(A)+P(B)-P(A\cap B)=0{,}7+0{,}5-0{,}3=\mathbf{0{,}9}$}

  \item $A$ et $B$ incompatibles, $P(A)=0{,}2$, $P(B)=0{,}2$. Calculer $P(A\cup B)$.

  \cor{$A$ et $B$ incompatibles $\Rightarrow P(A\cap B)=0$, donc $P(A\cup B)=0{,}2+0{,}2=\mathbf{0{,}4}$}

  \item $P(\bar{A})=0{,}2$, $P(B)=0{,}4$, $P(A\cap B)=0{,}2$. Calculer $P(A\cup B)$.

  \cor{$P(A)=1-P(\bar{A})=0{,}8$, donc $P(A\cup B)=0{,}8+0{,}4-0{,}2=\mathbf{1}$}
\end{enumerate}

\vfill

%------------------------------------------------------------------
\bex{Jeux à gratter}{%
Une société souhaite mettre sur le marché un nouveau jeu à gratter. La première case présente un soleil ou une lune ; la seconde case présente un cœur ou un pique. Sur \np{10000} tickets :
\begin{itemize}
  \item $10\,\%$ comportent un soleil ;
  \item $1\,\%$ des tickets avec un soleil comportent un cœur ;
  \item $0{,}5\,\%$ des tickets avec une lune comportent un cœur.
\end{itemize}
}

\textbf{1.} Compléter le tableau :

\begin{center}
\begin{tabular}{|c|c|c|c|}
\hline
\rowcolor{white} & Avec un soleil & Avec une lune & Total \\
\hline
\rowcolor{white} Avec un cœur  & \ans{10} & \ans{45} & \ans{55} \\
\hline
\rowcolor{white} Avec un pique & \ans{990} & \ans{\np{8955}} & \ans{\np{9945}} \\
\hline
\rowcolor{white} Total & \ans{\np{1000}} & \ans{\np{9000}} & \np{10000} \\
\hline
\end{tabular}
\end{center}

\cor{$10\,\%\times\np{10000}=\np{1000}$ soleils. $1\,\%\times\np{1000}=10$ (soleil+cœur). $0{,}5\,\%\times\np{9000}=45$ (lune+cœur).}

\textbf{2.} Calculer la probabilité que le ticket comporte un cœur.

\cor{$P(\text{cœur})=\dfrac{55}{\np{10000}}=\mathbf{0{,}0055}$}

\vfill

%------------------------------------------------------------------
\bex{QCM -- Machines A et B}{%
Deux machines A et B produisent une pièce. On a prélevé \np{3000} pièces de A et \np{2000} de B.
Pour A~: $2\,\%$ défaut C seul, $5\,\%$ défaut T seul, $1\,\%$ les deux défauts.
Pour B~: $3\,\%$ défaut C seul, $4\,\%$ défaut T seul, $2\,\%$ les deux défauts.
}

\begin{center}
\begin{tabular}{|c|c|c|c|c|c|}
\hline
\rowcolor{white} & C seul & T seul & C et T & Ni C ni T & Total \\
\hline
\rowcolor{white} Machine A & \ans{60} & \ans{150} & \ans{30} & \ans{\np{2760}} & \ans{\np{3000}} \\
\hline
\rowcolor{white} Machine B & \ans{60} & \ans{80} & \ans{40} & \ans{\np{1820}} & \ans{\np{2000}} \\
\hline
\rowcolor{white} Total & \ans{120} & \ans{230} & \ans{70} & \ans{\np{4580}} & \np{5000} \\
\hline
\end{tabular}
\end{center}

\begin{multicols}{2}
\noindent 1. La probabilité que la pièce vienne de A :
\begin{itemize}[label=$\square$]
  \item \ans{\ding{52}} a) $\dfrac{3}{5}$
  \item b) $\dfrac{2}{5}$
  \item c) Le tableau ne permet pas de conclure
\end{itemize}

2. La proba pour le défaut C seul :
\begin{itemize}[label=$\square$]
  \item \ans{\ding{52}} a) $0{,}024$
  \item b) $0{,}02$
  \item c) $0{,}03$
  \item d) $0{,}120$
\end{itemize}

\columnbreak

3. La proba pour le défaut T seul :
\begin{itemize}[label=$\square$]
  \item \ans{\ding{52}} a) $\dfrac{23}{500}$
  \item b) $\dfrac{3}{50}$
  \item c) $\dfrac{7}{500}$
  \item d) $\dfrac{1}{20}$
\end{itemize}

4. La proba pour au moins l'un des défauts :
\begin{itemize}[label=$\square$]
  \item a) $0{,}014$
  \item b) $0{,}06$
  \item c) $0{,}038$
  \item \ans{\ding{52}} d) $0{,}084$
\end{itemize}
\end{multicols}

\cor{\textit{Détail~:} $\frac{120}{5000}=0{,}024$ ; $\frac{230}{5000}=\frac{23}{500}$ ; $\frac{120+230+70}{5000}=\frac{420}{5000}=0{,}084$.}

\vfill

%------------------------------------------------------------------
\bpre{\textbf{Pièces défectueuses dans la production}\\[2pt]
Dans un atelier, deux machines A et B produisent des pièces. $40\,\%$ proviennent de A dont $5\,\%$ sont défectueuses. $2\,\%$ des pièces de B sont défectueuses.

\textbf{1.} Compléter le tableau ($\np{5000}$ pièces en tout) :

\begin{center}
\begin{tabular}{|c|c|c|c|}
\hline
\rowcolor{white} & Produites par A & Produites par B & Total \\
\hline
\rowcolor{white} Défectueuses & \ans{100} & \ans{60} & \ans{160} \\
\hline
\rowcolor{white} Non défectueuses & \ans{\np{1900}} & \ans{\np{2940}} & \ans{\np{4840}} \\
\hline
\rowcolor{white} Total & \ans{\np{2000}} & \ans{\np{3000}} & \np{5000} \\
\hline
\end{tabular}
\end{center}

\cor{$40\,\%\times\np{5000}=\np{2000}$ pièces de A ; $5\,\%\times\np{2000}=100$ défectueuses de A ; $2\,\%\times\np{3000}=60$ défectueuses de B.}

\textbf{2.} Calculer $P(A)$, $P(B)$ et $P(D)$ :

\begin{multicols}{3}
\cor{$P(A)=\dfrac{\np{2000}}{\np{5000}}=\mathbf{0{,}4}$}
\cor{$P(B)=\dfrac{\np{3000}}{\np{5000}}=\mathbf{0{,}6}$}
\cor{$P(D)=\dfrac{160}{\np{5000}}=\mathbf{0{,}032}$}
\end{multicols}

\textbf{3.} Définir $A\cap D$ et calculer sa probabilité.

\cor{$A\cap D$~: \og la pièce est de A \textbf{et} défectueuse\fg. \quad$P(A\cap D)=\dfrac{100}{\np{5000}}=0{,}02$}

\textbf{4.} Parmi les pièces défectueuses, calculer la probabilité qu'elle vienne de A.

\cor{Sur $160$ pièces défectueuses, $100$ viennent de A~:
$$P_D(A)=\frac{100}{160}=\frac{5}{8}=\mathbf{0{,}625}$$}

\textbf{5.} Calculer $\dfrac{P(A\cap D)}{P(D)}$ et comparer.

\cor{$\dfrac{P(A\cap D)}{P(D)}=\dfrac{0{,}02}{0{,}032}=\dfrac{5}{8}=0{,}625=P_D(A)$ \quad\checkmark\quad C'est la définition de la probabilité conditionnelle.}
}

\vfill

%------------------------------------------------------------------
\bpre{\textbf{Tirage dans une urne}\\[2pt]
Une urne contient 3 boules blanches (B) et 2 boules rouges (R). On tire successivement 2 boules sans remise.

\begin{multicols}{2}
\begin{center}
\begin{tikzpicture}[scale=0.55]
\fill (0,0) circle (1pt);
\node[circle,draw] (r1) at (2.5,2)  {R};
\node[circle,draw] (b1) at (2.5,-2) {B};
\node[circle,draw] (r21) at (5.5,3)  {R};
\node[circle,draw] (b21) at (5.5,1)  {B};
\node[circle,draw] (r22) at (5.5,-1) {R};
\node[circle,draw] (b22) at (5.5,-3) {B};
\draw (0,0) -- (r1); \draw (0,0) -- (b1);
\draw (r1) -- (r21); \draw (r1) -- (b21);
\draw (b1) -- (r22); \draw (b1) -- (b22);
% Probabilités
\node[above,blue!80!black,font=\small] at (1.2,1.2) {$\frac{2}{5}$};
\node[below,blue!80!black,font=\small] at (1.2,-1.2) {$\frac{3}{5}$};
\node[above,blue!80!black,font=\tiny] at (4,2.6) {$\frac{1}{4}$};
\node[below,blue!80!black,font=\tiny] at (4,1.4) {$\frac{3}{4}$};
\node[above,blue!80!black,font=\tiny] at (4,-1.4) {$\frac{2}{4}$};
\node[below,blue!80!black,font=\tiny] at (4,-2.6) {$\frac{2}{4}$};
\end{tikzpicture}
\end{center}

\begin{enumerate}
  \item Nombre de tirages distincts possibles.
  \cor{$5\times4=20$ tirages ordonnés (ou $\binom{5}{2}=10$ non ordonnés).}
  \item Probabilité d'obtenir deux rouges.
  \cor{$P(RR)=\frac{2}{5}\times\frac{1}{4}=\frac{1}{10}$}
  \item Probabilité d'obtenir deux blanches.
  \cor{$P(BB)=\frac{3}{5}\times\frac{2}{4}=\frac{3}{10}$}
  \item Probabilité d'obtenir deux couleurs différentes.
  \cor{$P=1-\frac{1}{10}-\frac{3}{10}=\frac{6}{10}=\frac{3}{5}$}
\end{enumerate}
\end{multicols}
}

\vfill

%------------------------------------------------------------------
\bpre{\textbf{Contrôle de production}\\[2pt]
Le contrôle de \np{10000} pièces a donné :

\begin{center}
\begin{tabular}{|c|c|c|c|}
\hline
& Défaut longueur & Pas de défaut longueur & Total \\
\hline
Défaut surface & 6 & 294 & 300 \\
\hline
Pas de défaut surface & 194 & \np{9506} & \np{9700} \\
\hline
Total & 200 & \np{9800} & \np{10000} \\
\hline
\end{tabular}
\end{center}

\textbf{1.} Calculer $P(L)$ et $P(S)$.

\begin{multicols}{2}
\cor{$P(L)=\dfrac{200}{\np{10000}}=\mathbf{0{,}02}$}
\cor{$P(S)=\dfrac{300}{\np{10000}}=\mathbf{0{,}03}$}
\end{multicols}

\textbf{2.} Calculer $P(S\cap L)$. Comparer $P(S\cap L)$ et $P(S)\times P(L)$.

\cor{$P(S\cap L)=\dfrac{6}{\np{10000}}=0{,}0006$\quad;\quad$P(S)\times P(L)=0{,}03\times0{,}02=0{,}0006$\\[4pt]
On a $P(S\cap L)=P(S)\times P(L)$ : $S$ et $L$ sont \textbf{indépendants}.}
}

%%%%%%%%%%%%%%%%%%%%%%%%%%
\newpage
%%%%%%%%%%%%%%%%%%%%%%%%%%

\begin{center}
\begin{minipage}{14cm}
\begin{bclogo}[couleur = white,logo=\bcinfo, arrondi = 0.1]{~Introduction~:}
~~
On fait remonter à la correspondance de $1654$ entre Pascal et Fermat l'acte de naissance du calcul des probabilités.

Après trois siècles de recherche, le calcul des probabilités a pu fournir les bases théoriques nécessaires au développement de la statistique.

Le but des probabilités est de rationaliser le hasard : quelles sont les chances d'obtenir un résultat suite à une expérience aléatoire ?
~~
\end{bclogo}
\end{minipage}
\end{center}

%%%%%%%%%%%%%%%%%%%%%%%%%%
\begin{mybox}
\begin{center}\textbf{1. Vocabulaire des événements}\end{center}
\end{mybox}

\begin{mybox}
\begin{center}\textbf{1.1. Vocabulaire}\end{center}
\end{mybox}

\begin{itemize}[label=\ding{228}]
  \item Chaque résultat possible d'une expérience aléatoire est une \underline{éventualité}.
  \item L'ensemble des éventualités est l'\underline{univers}, noté $\Omega$.
  \item Un \underline{événement} est une partie de $\Omega$.
  \item Un événement à une seule éventualité est un \underline{événement élémentaire}.
  \item L'\underline{événement impossible} $\varnothing$ ne contient aucune éventualité.
  \item L'\underline{événement certain} est $\Omega$.
  \item L'\underline{événement contraire} de $A$ est $\bar{A}$ : les éléments de $\Omega$ qui ne sont pas dans $A$. On a $A\cup\bar{A}=\Omega$ et $A\cap\bar{A}=\varnothing$.
\end{itemize}

\cor{%
\textbf{Exemple~:} On lance un dé à 6 faces. $\Omega=\{1,2,3,4,5,6\}$.\\
$A=\text{\og nombre pair\fg}=\{2,4,6\}$\quad;\quad$\bar{A}=\{1,3,5\}$ (impair).\\
\og obtenir 7\fg{} est \textbf{impossible} ($\varnothing$)\quad;\quad$\Omega$ est \textbf{certain}.}

Dans toute la suite, $A$ et $B$ désignent deux événements de l'univers $\Omega$.

\newpage

%%%%%%%%%%%%%%%%%%%%%%%%%%
\begin{mybox}
\begin{center}\textbf{1.2. Intersection et réunion des événements}\end{center}
\end{mybox}

\fbox{\begin{minipage}{16.6cm}
\begin{itemize}[label=\ding{228}]
  \item \underline{Intersection} $A\cap B$ (« $A$ et $B$ ») : éventualités appartenant à $A$ \textbf{et} à $B$.
  \item \underline{Réunion} $A\cup B$ (« $A$ ou $B$ ») : éventualités appartenant à $A$ \textbf{ou} à $B$.
\end{itemize}
\begin{center}
  Si $A\cap B=\varnothing$, on dit que $A$ et $B$ sont \underline{disjoints} ou \underline{incompatibles}.
\end{center}
\end{minipage}}

\cor{%
\textbf{Exemple~:} $\Omega=\{1,2,3,4,5,6\}$, $A=\{2,4,6\}$ (pair), $B=\{3,6\}$ (multiple de 3).\\
$A\cap B=\{6\}$\quad;\quad$A\cup B=\{2,3,4,6\}$.\\[4pt]
\begin{center}
\begin{tikzpicture}[scale=0.7]
\draw[thick] (-2.5,-1.7) rectangle (2.5,1.7);
\node at (-2.1,1.4) {$\Omega$};
\draw[fill=blue!20,opacity=0.5] (-0.7,0) circle (1.2cm);
\node at (-1.2,0) {$A$};
\draw[fill=green!20,opacity=0.5] (0.7,0) circle (1.2cm);
\node at (1.2,0) {$B$};
\node at (0,0) {\small$A\cap B$};
\node[below] at (0,-1.4) {\small$A\cup B$ = zones bleue et verte};
\end{tikzpicture}
\end{center}}

%%%%%%%%%%%%%%%%%%%%%%%%%%
\begin{mybox}
\begin{center}\textbf{1.3. Représentation des événements}\end{center}
\end{mybox}

\textbf{Diagrammes de Venn}

\begin{center}
\begin{multicols}{4}
\begin{tikzpicture}[scale=0.45]
\newcommand{\cA}{(0,0.67) circle (1cm)}
\newcommand{\cB}{(0,-0.67) circle (1cm)}
\fill[blue!20,even odd rule] \cA \cB;
\draw[blue] \cA; \draw[blue] \cB;
\node[above] at (0,-2.3) {\small\ans{$A\cup B$}};
\end{tikzpicture}

\begin{tikzpicture}[scale=0.45]
\newcommand{\cA}{(0,0.67) circle (1cm)}
\newcommand{\cB}{(0,-0.67) circle (1cm)}
\draw[blue] \cA; \draw[blue] \cB;
\begin{scope}\clip \cA;\fill[blue!20] \cB;\end{scope}
\node[above] at (0,-2.3) {\small\ans{$A\cap B$}};
\end{tikzpicture}

\begin{tikzpicture}[scale=0.45]
\newcommand{\cA}{(0,0.67) circle (1cm)}
\newcommand{\cB}{(0,-0.67) circle (1cm)}
\begin{scope}\clip(-1.1,0)rectangle(1.1,1.8);\fill[blue!20]\cA;\draw[blue]\cA;\end{scope}
\begin{scope}\clip(-1.1,-1.8)rectangle(1.1,0);\fill[blue!20]\cB;\draw[blue]\cB;\end{scope}
\node[above] at (0,-2.3) {\small\ans{$A\cup B$, $A\cap B=\varnothing$}};
\end{tikzpicture}

\begin{tikzpicture}[scale=0.45]
\newcommand{\cA}{(0,0.67) circle (1cm)}
\newcommand{\cB}{(0,-0.67) circle (1cm)}
\fill[blue!20] \cB;
\draw[blue,fill=white] \cA;
\draw[blue] \cB;
\node[above] at (0,-2.3) {\small\ans{$B\cap\bar{A}$}};
\end{tikzpicture}
\end{multicols}
\end{center}

\vfill

\textbf{Tableau}~~On jette deux dés à 4 faces et on calcule le \textbf{produit} :

\begin{center}
\begin{tabular}{|c|c|c|c|c|}
\hline
$\times$ & \textbf{1} & \textbf{2} & \textbf{3} & \textbf{4} \\
\hline
\textbf{1} & \ans{1} & \ans{2} & \ans{3} & \ans{4} \\
\hline
\textbf{2} & \ans{2} & \ans{4} & \ans{6} & \ans{8} \\
\hline
\textbf{3} & \ans{3} & \ans{6} & \ans{9} & \ans{12} \\
\hline
\textbf{4} & \ans{4} & \ans{8} & \ans{12} & \ans{16} \\
\hline
\end{tabular}
\end{center}

\vfill

\textbf{Arbre}~~On lance une pièce 3 fois ($P=$ pile, $F=$ face) :

\begin{center}
\begin{tikzpicture}[scale=0.5]
\fill(0,0)circle(1pt);
\foreach \n/\y in {P/4,F/-4}
  \node[circle,draw](l1\n) at(3,\y){\n};
\foreach \n/\y in {P/6,F/2}
  \node[circle,draw](l2P\n) at(6,\y){\n};
\foreach \n/\y in {P/-2,F/-6}
  \node[circle,draw](l2F\n) at(6,\y){\n};
\foreach \n/\y in {P/7,F/5}
  \node[circle,draw](l3PP\n) at(9,\y){\n};
\foreach \n/\y in {P/3,F/1}
  \node[circle,draw](l3PF\n) at(9,\y){\n};
\foreach \n/\y in {P/-1,F/-3}
  \node[circle,draw](l3FP\n) at(9,\y){\n};
\foreach \n/\y in {P/-5,F/-7}
  \node[circle,draw](l3FF\n) at(9,\y){\n};
\draw(0,0)--(l1P); \draw(0,0)--(l1F);
\draw(l1P)--(l2PP); \draw(l1P)--(l2PF);
\draw(l1F)--(l2FP); \draw(l1F)--(l2FF);
\draw(l2PP)--(l3PPP); \draw(l2PP)--(l3PPF);
\draw(l2PF)--(l3PFP); \draw(l2PF)--(l3PFF);
\draw(l2FP)--(l3FPP); \draw(l2FP)--(l3FPF);
\draw(l2FF)--(l3FFP); \draw(l2FF)--(l3FFF);
\node[above,blue!80!black,font=\tiny] at(1.5,2.3){$\frac{1}{2}$};
\node[below,blue!80!black,font=\tiny] at(1.5,-2.3){$\frac{1}{2}$};
\end{tikzpicture}
\end{center}

\newpage

%%%%%%%%%%%%%%%%%%%%%%%%%%
\begin{mybox}
\begin{center}\textbf{2. Calcul de probabilités}\end{center}
\end{mybox}

\bdef{%
\textcolor{blue!80!black}{%
Dans un univers $\Omega$ fini à $n$ éventualités \textbf{équiprobables}, la probabilité d'un événement $A$ est~:
$$\boxed{P(A)=\frac{\text{Card}(A)}{\text{Card}(\Omega)}=\frac{\text{nombre d'éventualités favorables}}{\text{nombre total d'éventualités}}}$$
}}

\rem{Dans un exercice, on reconnaît l'équiprobabilité aux expressions : dé \underline{non pipé}, boules \underline{indiscernables} au toucher, tirage au \underline{hasard}\dots}

\vfill

\bprop{}{%
Soit $A$ et $B$ deux événements :
\begin{itemize}[label=\ding{217}]
  \item $P(\varnothing)=0$ \quad;\quad $P(\Omega)=1$ \quad;\quad $0\leq P(A)\leq 1$
  \item $P(\bar{A})=1-P(A)$
  \item $P(A\cup B)=P(A)+P(B)-P(A\cap B)$
\end{itemize}
}

\vfill

%------------------------------------------------------------------
\bex{}{On considère $E=\{1,2,\ldots,20\}$. On choisit un nombre au hasard.}

\begin{enumerate}
  \item $A$ : \og le nombre est multiple de $3$ \fg

  \cor{$A=\{3,6,9,12,15,18\}$ \quad Card$(A)=6$}

  \item $B$ : \og le nombre est multiple de $2$ \fg

  \cor{$B=\{2,4,6,8,10,12,14,16,18,20\}$ \quad Card$(B)=10$}

  \item Calculer $P(A)$, $P(\bar{A})$, $P(B)$, $P(A\cap B)$ et $P(A\cup B)$.

  \cor{%
  $P(A)=\dfrac{6}{20}=\dfrac{3}{10}=0{,}3$\quad;\quad
  $P(\bar{A})=1-\dfrac{3}{10}=\dfrac{7}{10}=0{,}7$\\[4pt]
  $P(B)=\dfrac{10}{20}=\dfrac{1}{2}=0{,}5$\\[4pt]
  $A\cap B=\{6,12,18\}$ (multiples de 6) \quad$\Rightarrow$\quad$P(A\cap B)=\dfrac{3}{20}=0{,}15$\\[4pt]
  $P(A\cup B)=\dfrac{3}{10}+\dfrac{1}{2}-\dfrac{3}{20}=\dfrac{6+10-3}{20}=\dfrac{13}{20}=0{,}65$}
\end{enumerate}

\vfill
\newpage

%%%%%%%%%%%%%%%%%%%%%%%%%%
\begin{mybox}
\begin{center}\textbf{3. Probabilités conditionnelles}\end{center}
\end{mybox}

\begin{mybox}
\begin{center}\textbf{3.1. Définition}\end{center}
\end{mybox}

\bdef{%
\textcolor{blue!80!black}{%
Soient $A$ et $B$ deux événements avec $P(B)\neq 0$. La \textbf{probabilité conditionnelle de $A$ sachant $B$} est~:
$$\boxed{P_B(A)=P(A|B)=\frac{P(A\cap B)}{P(B)}}$$
On note aussi $P(A/B)$.}}

\vfill

%------------------------------------------------------------------
\bex{}{On lance un dé à 6 faces équilibré. $B$~: \og numéro pair\fg. $A$~: \og numéro multiple de $3$\fg.

Calculer $P_B(A)$ de deux manières.}

\begin{multicols}{2}
\cor{\textbf{Méthode 1 (univers réduit)}\\
$B=\{2,4,6\}$. Parmi ces 3 valeurs, seul $6$ est multiple de $3$.\\[4pt]
$P_B(A)=\dfrac{1}{3}$}
\cor{\textbf{Méthode 2 (formule)}\\
$P(B)=\frac{3}{6}=\frac{1}{2}$\quad;\quad$A\cap B=\{6\}$, $P(A\cap B)=\frac{1}{6}$\\[4pt]
$P_B(A)=\dfrac{1/6}{1/2}=\dfrac{1}{3}$}
\end{multicols}

\vfill

\bprop{Formule des probabilités composées}{%
Pour $P(B)\neq0$ et $P(A)\neq0$~:
$$\boxed{P(A\cap B)=P(B)\times P_B(A)=P(A)\times P_A(B)}$$}

\textit{Démonstration~:} $P_B(A)=\dfrac{P(A\cap B)}{P(B)}\Rightarrow P(A\cap B)=P(B)\,P_B(A)$ et de même pour $P_A(B)$. \hfill$\square$

\vfill

\bprop{$P_B$ est une probabilité}{%
Si $P(B)\neq0$, $P_B$ vérifie les mêmes propriétés que $P$~:
\begin{itemize}
  \item $P_B(\varnothing)=0$ \quad;\quad $P_B(\Omega)=P_B(B)=1$ \quad;\quad $0\leq P_B(A)\leq1$
  \item $P_B(\bar{A})=1-P_B(A)$
  \item $P_B(A\cup C)=P_B(A)+P_B(C)-P_B(A\cap C)$
\end{itemize}}

\vfill

\bprop{Formule des probabilités totales}{%
Si $\{B,\bar{B}\}$ est une partition de $\Omega$~:
$$\boxed{P(A)=P(B)\times P_B(A)+P(\bar{B})\times P_{\bar{B}}(A)}$$
Plus généralement, si $\{B_1,\ldots,B_n\}$ est une partition de $\Omega$~:
$$P(A)=\sum_{i=1}^{n}P(B_i)\times P_{B_i}(A)$$}

\textit{Démonstration~:} $P(A)=P(A\cap B)+P(A\cap\bar{B})=P(B)\,P_B(A)+P(\bar{B})\,P_{\bar{B}}(A)$. \hfill$\square$

\newpage

%%%%%%%%%%%%%%%%%%%%%%%%%%
\begin{mybox}
\begin{center}\textbf{3.2. Exemple de calcul par différentes méthodes}\end{center}
\end{mybox}

\bex{}{%
Dans un atelier, $M_1$ fournit $60\,\%$ de la production (dont $6{,}3\,\%$ défectueuses) et $M_2$ le reste (dont $4\,\%$ défectueuses). On tire une pièce au hasard.}

\begin{enumerate}
\item \textbf{Formules de probabilités conditionnelles}

\begin{enumerate}
  \item Probabilité d'une pièce défectueuse sachant qu'elle vient de $M_1$ ?

  \ans{$P_{M_1}(D)=0{,}063$}

  \item Probabilité d'une pièce défectueuse sachant qu'elle vient de $M_2$ ?

  \ans{$P_{M_2}(D)=0{,}04$}

  \item Probabilité de prélever une pièce défectueuse ?

  \cor{$P(D)=P(M_1)\times P_{M_1}(D)+P(M_2)\times P_{M_2}(D)=0{,}6\times0{,}063+0{,}4\times0{,}04=0{,}0378+0{,}016=\mathbf{0{,}0538}$}
\end{enumerate}

\item \textbf{Tableau} (production de \np{10000} pièces)

\begin{center}
\renewcommand{\arraystretch}{1.3}
\begin{tabular}{|l|c|c|c|}
\hline
& Pièces de $M_1$ & Pièces de $M_2$ & Total \\
\hline
Défectueuses & \ans{378} & \ans{160} & \ans{538} \\
\hline
Conformes & \ans{\np{5622}} & \ans{\np{3840}} & \ans{\np{9462}} \\
\hline
Total & \ans{\np{6000}} & \ans{\np{4000}} & \np{10000} \\
\hline
\end{tabular}
\end{center}

\cor{$P(D)=\dfrac{538}{\np{10000}}=\mathbf{0{,}0538}$ \quad (même résultat \checkmark)}

\item \textbf{Arbre des probabilités conditionnelles}

\cor{%
\begin{tikzpicture}[scale=0.9,every node/.style={font=\small}]
\fill(0,0)circle(1pt);
\node(m1) at(3,1.5){$M_1$}; \node(m2) at(3,-1.5){$M_2$};
\node(d1) at(6,2.5){$D$}; \node(nd1) at(6,0.5){$\bar{D}$};
\node(d2) at(6,-0.5){$D$}; \node(nd2) at(6,-2.5){$\bar{D}$};
\draw(0,0)--(m1)node[midway,above]{\textcolor{blue!80!black}{$0{,}6$}};
\draw(0,0)--(m2)node[midway,below]{\textcolor{blue!80!black}{$0{,}4$}};
\draw(m1)--(d1)node[midway,above]{\textcolor{blue!80!black}{$0{,}063$}};
\draw(m1)--(nd1)node[midway,below]{\textcolor{blue!80!black}{$0{,}937$}};
\draw(m2)--(d2)node[midway,above]{\textcolor{blue!80!black}{$0{,}04$}};
\draw(m2)--(nd2)node[midway,below]{\textcolor{blue!80!black}{$0{,}96$}};
\node[right,blue!80!black,font=\scriptsize] at(6.3,2.5){$P(M_1\cap D)=0{,}6\times0{,}063=0{,}0378$};
\node[right,blue!80!black,font=\scriptsize] at(6.3,0.5){$P(M_1\cap\bar{D})=0{,}6\times0{,}937=0{,}5622$};
\node[right,blue!80!black,font=\scriptsize] at(6.3,-0.5){$P(M_2\cap D)=0{,}4\times0{,}04=0{,}016$};
\node[right,blue!80!black,font=\scriptsize] at(6.3,-2.5){$P(M_2\cap\bar{D})=0{,}4\times0{,}96=0{,}384$};
\end{tikzpicture}

$P(D)=0{,}0378+0{,}016=\mathbf{0{,}0538}$ \checkmark}
\end{enumerate}

\newpage

%%%%%%%%%%%%%%%%%%%%%%%%%%
\begin{mybox}
\begin{center}\textbf{3.3. Événements indépendants}\end{center}
\end{mybox}

\bdef{%
\textcolor{blue!80!black}{%
$A$ et $B$ sont \textbf{indépendants} si et seulement si~:
$$\boxed{P(A\cap B)=P(A)\times P(B)}$$
Conséquence~: si $P(B)\neq0$, $A$ et $B$ indépendants $\Leftrightarrow P_B(A)=P(A)$ (connaître $B$ ne modifie pas la probabilité de $A$).}}

\rem{Indépendance $\neq$ incompatibilité. Deux événements incompatibles de probabilité non nulle ne sont pas indépendants.}

\vfill

%------------------------------------------------------------------
\bex{}{On tire au hasard une carte d'un jeu de 32 cartes. $A=$~tirer un as, $B=$~tirer un cœur, $C=$~tirer un as rouge.}

\textbf{Indépendance de $A$ et $B$ :}
\cor{%
$P(A)=\dfrac{4}{32}=\dfrac{1}{8}$\quad;\quad
$P(B)=\dfrac{8}{32}=\dfrac{1}{4}$\quad;\quad
$P(A\cap B)=\dfrac{1}{32}$ (as de cœur)\\[4pt]
$P(A)\times P(B)=\dfrac{1}{8}\times\dfrac{1}{4}=\dfrac{1}{32}=P(A\cap B)$ \quad$\Rightarrow$ $A$ et $B$ sont \textbf{indépendants}.}

\textbf{Indépendance de $B$ et $C$ :}
\cor{%
$P(C)=\dfrac{2}{32}=\dfrac{1}{16}$ (as rouge)\quad;\quad
$P(B\cap C)=\dfrac{1}{32}$ (as de cœur)\\[4pt]
$P(B)\times P(C)=\dfrac{1}{4}\times\dfrac{1}{16}=\dfrac{1}{64}\neq\dfrac{1}{32}=P(B\cap C)$ \quad$\Rightarrow$ $B$ et $C$ ne sont \textbf{pas indépendants}.}

\vfill

\bprop{}{Si $A$ et $B$ sont indépendants, alors $A$ et $\bar{B}$ ; $\bar{A}$ et $B$ ; $\bar{A}$ et $\bar{B}$ sont aussi indépendants.}

\newpage

%------------------------------------------------------------------
\bex{}{%
\begin{center}\textit{Les résultats seront arrondis à $10^{-5}$.}\end{center}
Un atelier assemble un moteur. $D$~: \og la machine est défaillante\fg. $A$~: \og l'alerte est donnée\fg.\\
$P(D)=0{,}001$\,;\; l'alarme se déclenche dans 99 cas sur 100 si défaillante\,;\; dans 5 cas sur 1000 si non défaillante.}

\begin{enumerate}
\item Donner $P(D)$, $P_D(A)$ et $P_{\bar{D}}(A)$.

\begin{multicols}{3}
\cor{$P(D)=\mathbf{0{,}001}$}
\cor{$P_D(A)=\mathbf{0{,}99}$}
\cor{$P_{\bar{D}}(A)=\dfrac{5}{1000}=\mathbf{0{,}005}$}
\end{multicols}

\item
\begin{enumerate}
  \item Calculer $P(A\cap D)$ et $P(A\cap\bar{D})$.

  \cor{$P(A\cap D)=0{,}001\times0{,}99=\mathbf{0{,}00099}$\\[2pt]
  $P(A\cap\bar{D})=0{,}999\times0{,}005=\mathbf{0{,}004995}$}

  \item En déduire $P(A)$.

  \cor{$P(A)=0{,}00099+0{,}004995=\mathbf{0{,}005985}$}
\end{enumerate}

\item L'alerte est donnée. Calculer $P_A(\bar{D})$ (probabilité de fausse alerte) à $10^{-2}$.

\cor{$P_A(\bar{D})=\dfrac{P(A\cap\bar{D})}{P(A)}=\dfrac{0{,}004995}{0{,}005985}\approx\mathbf{0{,}83}$\\[4pt]
\textbf{Interprétation~:} Quand l'alarme sonne, il y a $83\,\%$ de chances que ce soit une fausse alerte.}
\end{enumerate}

\vfill

%------------------------------------------------------------------
\bex{}{%
$C$~: \og la calculatrice a un défaut de clavier\fg\,;\; $A$~: \og la calculatrice a un défaut d'affichage\fg.\\
$P(C)=0{,}04$\,;\; $P_C(A)=0{,}03$\,;\; $P_{\bar{C}}(\bar{A})=0{,}94$. Probabilités arrondies au millième.}

\begin{enumerate}
\item Identifier les probabilités de l'énoncé.

\cor{$P_{\bar{C}}(\bar{A})=0{,}94$\quad;\quad
$P_{\bar{C}}(A)=1-0{,}94=0{,}06$\quad;\quad
$P(C)=0{,}04$, $P(\bar{C})=0{,}96$}

\item Construire un arbre pondéré.

\cor{%
\begin{tikzpicture}[scale=0.85,every node/.style={font=\small}]
\fill(0,0)circle(1pt);
\node(c) at(3,1.5){$C$}; \node(cb) at(3,-1.5){$\bar{C}$};
\node(a1) at(6,2.3){$A$}; \node(ab1) at(6,0.7){$\bar{A}$};
\node(a2) at(6,-0.7){$A$}; \node(ab2) at(6,-2.3){$\bar{A}$};
\draw(0,0)--(c)node[midway,above]{\textcolor{blue!80!black}{$0{,}04$}};
\draw(0,0)--(cb)node[midway,below]{\textcolor{blue!80!black}{$0{,}96$}};
\draw(c)--(a1)node[midway,above]{\textcolor{blue!80!black}{$0{,}03$}};
\draw(c)--(ab1)node[midway,below]{\textcolor{blue!80!black}{$0{,}97$}};
\draw(cb)--(a2)node[midway,above]{\textcolor{blue!80!black}{$0{,}06$}};
\draw(cb)--(ab2)node[midway,below]{\textcolor{blue!80!black}{$0{,}94$}};
\node[right,blue!80!black,font=\scriptsize] at(6.2,2.3){$0{,}04\times0{,}03=0{,}001$};
\node[right,blue!80!black,font=\scriptsize] at(6.2,0.7){$0{,}04\times0{,}97=0{,}039$};
\node[right,blue!80!black,font=\scriptsize] at(6.2,-0.7){$0{,}96\times0{,}06=0{,}058$};
\node[right,blue!80!black,font=\scriptsize] at(6.2,-2.3){$0{,}96\times0{,}94=0{,}902$};
\end{tikzpicture}}

\item
\begin{enumerate}
  \item Calculer la probabilité que la calculatrice présente les \textbf{deux} défauts.

  \cor{$P(A\cap C)=0{,}04\times0{,}03=\mathbf{0{,}001}$}

  \item Calculer la probabilité d'un défaut d'affichage sans défaut de clavier.

  \cor{$P(A\cap\bar{C})=0{,}96\times0{,}06=\mathbf{0{,}058}$}

  \item En déduire $P(A)$.

  \cor{$P(A)=0{,}001+0{,}058=\mathbf{0{,}059}$}

  \item Montrer que $P(\bar{A}\cap\bar{C})\approx0{,}902$.

  \cor{$P(\bar{A}\cap\bar{C})=P(\bar{C})\times P_{\bar{C}}(\bar{A})=0{,}96\times0{,}94=0{,}9024\approx\mathbf{0{,}902}$ \checkmark}
\end{enumerate}
\end{enumerate}

\newpage

%------------------------------------------------------------------
\bex{}{%
Machine A : \np{1500} disques/jour, machine B : \np{3000} disques/jour.\\
$P_A(D)=0{,}02$ (défaut sachant A)\,;\; $P_B(D)=0{,}035$ (défaut sachant B).}

\begin{enumerate}
\item Calculer $P(A)$, $P(B)$ et $P(D)$.

\cor{$P(A)=\dfrac{\np{1500}}{\np{4500}}=\dfrac{1}{3}\approx0{,}333$\quad;\quad
$P(B)=\dfrac{\np{3000}}{\np{4500}}=\dfrac{2}{3}\approx0{,}667$\\[4pt]
$P(D)=\dfrac{1}{3}\times0{,}02+\dfrac{2}{3}\times0{,}035=\dfrac{0{,}02+0{,}07}{3}=\dfrac{0{,}09}{3}=\mathbf{0{,}03}$}

\item Le disque a un défaut. Quelle est la probabilité qu'il vienne de A ? de B ?

\cor{$P_D(A)=\dfrac{P(A\cap D)}{P(D)}=\dfrac{(1/3)\times0{,}02}{0{,}03}=\dfrac{0{,}02}{0{,}09}=\dfrac{2}{9}\approx\mathbf{0{,}222}$\\[4pt]
$P_D(B)=\dfrac{P(B\cap D)}{P(D)}=\dfrac{(2/3)\times0{,}035}{0{,}03}=\dfrac{0{,}07}{0{,}09}=\dfrac{7}{9}\approx\mathbf{0{,}778}$\\[4pt]
Vérification~: $P_D(A)+P_D(B)=\frac{2}{9}+\frac{7}{9}=1$ \checkmark}
\end{enumerate}

\vfill

%------------------------------------------------------------------
\bex{}{%
Tableau d'élèves utilisant l'application Toktok (probabilités au millième)~:

\begin{center}
\begin{tabular}{|c|c|c|c|}
\hline
& Filles & Garçons & Total \\
\hline
Utilisent Toktok & 148 & 171 & 319 \\
\hline
N'utilisent pas Toktok & 81 & 50 & 131 \\
\hline
Total & 229 & 221 & 450 \\
\hline
\end{tabular}
\end{center}
$G$~: \og garçon\fg\,;\; $T$~: \og utilise Toktok\fg.}

\begin{enumerate}
\item Calculer $P(G)$.

\cor{$P(G)=\dfrac{221}{450}\approx\mathbf{0{,}491}$}

\item Montrer que $P(G\cap T)=0{,}38$.

\cor{$P(G\cap T)=\dfrac{171}{450}=0{,}38$ \checkmark}

\item Calculer la probabilité de prélever la fiche d'une fille sachant que l'élève n'utilise pas Toktok.

\cor{$P_{\bar{T}}(\bar{G})=\dfrac{P(\bar{G}\cap\bar{T})}{P(\bar{T})}=\dfrac{81/450}{131/450}=\dfrac{81}{131}\approx\mathbf{0{,}618}$}

\item Calculer $P_T(G)$ et interpréter.

\cor{$P_T(G)=\dfrac{171/450}{319/450}=\dfrac{171}{319}\approx\mathbf{0{,}536}$\\[4pt]
\textbf{Interprétation~:} Parmi les élèves utilisant Toktok, environ $53{,}6\,\%$ sont des garçons.}
\end{enumerate}

\newpage

%------------------------------------------------------------------
\bex{}{%
$A$~: \og comprimé conforme\fg\,;\; $B$~: \og comprimé refusé\fg.\\
$P(A)=0{,}98$\,;\; un comprimé conforme est \emph{accepté} avec prob. $0{,}98$\,;\; un non-conforme est refusé avec prob. $0{,}99$.}

\begin{enumerate}
\item Calculer $P_A(B)$, $P(B\cap A)$ et $P(B\cap\bar{A})$.

\cor{$P_A(B)=1-0{,}98=\mathbf{0{,}02}$ (conforme mais refusé)\\[4pt]
$P(B\cap A)=P(A)\times P_A(B)=0{,}98\times0{,}02=\mathbf{0{,}0196}$\\[4pt]
$P(\bar{A})=0{,}02$\quad;\quad$P_{\bar{A}}(B)=0{,}99$\quad;\quad
$P(B\cap\bar{A})=0{,}02\times0{,}99=\mathbf{0{,}0198}$}

\item
\begin{enumerate}
  \item La probabilité qu'un comprimé soit refusé~:

  \cor{$P(B)=P(B\cap A)+P(B\cap\bar{A})=0{,}0196+0{,}0198=\mathbf{0{,}0394}$}

  \item La probabilité qu'un comprimé refusé soit conforme (à $10^{-4}$)~:

  \cor{$P_B(A)=\dfrac{P(A\cap B)}{P(B)}=\dfrac{0{,}0196}{0{,}0394}\approx\mathbf{0{,}4975}$}
\end{enumerate}
\end{enumerate}

\vfill

%------------------------------------------------------------------
\bex{}{%
Un commercial visite deux clients. $P(A)=0{,}25$ ($A$~: le $1^{\text{er}}$ achète)\,;\; $P_A(B)=0{,}4$\,;\; $P_{\bar{A}}(B)=0{,}25$ ($B$~: le $2^{\text{e}}$ achète).}

\textbf{1.} Arbre pondéré :

\cor{%
\begin{tikzpicture}[scale=0.85,every node/.style={font=\small}]
\fill(0,0)circle(1pt);
\node(a) at(3,1.5){$A$}; \node(ab) at(3,-1.5){$\bar{A}$};
\node(b1) at(6,2.3){$B$}; \node(bb1) at(6,0.7){$\bar{B}$};
\node(b2) at(6,-0.7){$B$}; \node(bb2) at(6,-2.3){$\bar{B}$};
\draw(0,0)--(a)node[midway,above]{\textcolor{blue!80!black}{$0{,}25$}};
\draw(0,0)--(ab)node[midway,below]{\textcolor{blue!80!black}{$0{,}75$}};
\draw(a)--(b1)node[midway,above]{\textcolor{blue!80!black}{$0{,}4$}};
\draw(a)--(bb1)node[midway,below]{\textcolor{blue!80!black}{$0{,}6$}};
\draw(ab)--(b2)node[midway,above]{\textcolor{blue!80!black}{$0{,}25$}};
\draw(ab)--(bb2)node[midway,below]{\textcolor{blue!80!black}{$0{,}75$}};
\node[right,blue!80!black,font=\scriptsize] at(6.2,2.3){$P(A\cap B)=0{,}25\times0{,}4=0{,}10$};
\node[right,blue!80!black,font=\scriptsize] at(6.2,0.7){$P(A\cap\bar{B})=0{,}25\times0{,}6=0{,}15$};
\node[right,blue!80!black,font=\scriptsize] at(6.2,-0.7){$P(\bar{A}\cap B)=0{,}75\times0{,}25=0{,}1875$};
\node[right,blue!80!black,font=\scriptsize] at(6.2,-2.3){$P(\bar{A}\cap\bar{B})=0{,}75\times0{,}75=0{,}5625$};
\end{tikzpicture}}

\textbf{2.} Calculer $P(B)$.

\cor{$P(B)=P(A\cap B)+P(\bar{A}\cap B)=0{,}10+0{,}1875=\mathbf{0{,}2875}$}

\textbf{3.} Probabilité qu'un seul client achète.

\cor{$P(\text{un seul})=P(A\cap\bar{B})+P(\bar{A}\cap B)=0{,}15+0{,}1875=\mathbf{0{,}3375}$}

\end{document}
